\section{OpenLineage: execution history without automatic scientific acceptance}
\label{sec:exec24-openlineage-reuse}
A bank comparing collections policies needs to know which data, model,
assumptions and decision rule produced each result. That need makes a common
lineage vocabulary useful, but also creates a risk: a completed job can be
mistaken for a justified causal answer. OpenLineage supplies run, job and
dataset events with extensible facets. At the selected 7 October 2026 commit,
the inspected core schema identifies itself as version~2.0.2 and the Python
package declares version~1.54.0. These pinned sources support a concrete
reference choice for the proposal. They do not establish that an existing
server preserves every facet, that another workspace has adopted the
contract, or that any emitted claim is scientifically true.
\citep{exec24_openlineage_source}

\subsection{A valid event is not a complete analysis history}
The core schema requires event time, producer URI and schema URI. A run
event additionally requires run and job objects; the run requires a UUID,
and the job requires namespace and name. The description of event type
calls for a START and one terminal COMPLETE, ABORT or FAIL per run, but
event type is not in that object's required-field list. Input and output
datasets are optional too. These are distinct facts: descriptive protocol
requirements are not all enforced by single-document JSON validation.
\citep{exec24_openlineage_source}

Local authored fixtures validate the acquired declarative schema using an
installed JSON Schema validator with format checking and only internal
references. A run record without event type passes; a COMPLETE record without
inputs or outputs passes; an invalid UUID fails. Two identical COMPLETE
records each pass although the pair has no START. These results concern the
schema, not native client construction, server ingestion or stream validation.
The document validator cannot establish a unique, complete lifecycle by
examining one event at a time. A product adapter therefore needs a separate
history rule with attempt identity, permitted state transitions, duplicate
handling and reconciliation of missing terminal records.
\citep{exec24_openlineage_reference}

Dataset namespace and name identify a dataset, but the selected base schema
does not require a content hash or immutable snapshot. Two authored payloads
with the same namespace/name and different hashes show why this matters.
The proposed causal facet should bind the exact data snapshot, availability
clock, eligible population, event dictionary, target, graph or model,
identification assumptions, nuisance fits, split assignment and uncertainty
method. Large or restricted artifacts should be referenced by authorized
identifiers and digests, not copied indiscriminately into every event.
\citep{exec24_openlineage_source,exec24_openlineage_reference}

The base facet requires producer and schema URIs and permits additional
properties. An authored facet asserting that a result is identified passes
the selected base schema despite also declaring an unsupported claim. The
test does not dereference its custom URI: that URI is a field, not proof that
an additional schema or its scientific claims were checked. A production
consumer must validate an approved custom-facet schema and separately review
the evidence supporting the assertion. Serialization, schema conformance,
causal identification and bank acceptance are separate states.
\citep{exec24_openlineage_source,exec24_openlineage_reference}

\subsection{Delivery guarantees depend on the concrete client and transport}
The Java client documents a normal return after emission regardless of backend
exceptions. The inspected implementation delegates through a metric timer
without an exception guard. Its concrete HTTP transport performs a synchronous
POST and raises on response codes from 400 through 599; its response exception
is a client exception. The comment is therefore not a sufficient universal
no-throw contract. The circuit-breaker branch can also return without emitting.
Even an incremented emission-complete counter is not a scientific acceptance
receipt. The wrapper must distinguish suppression, attempted delivery,
acknowledged delivery, persisted lineage and accepted analysis.
\citep{exec24_openlineage_source}

The Java duration overload introduces another boundary. The abstract
transport's default implementation accepts a duration but simply calls the
ordinary emit method. The inspected HTTP transport does not override that
overload; it configures connection, socket and response timeouts separately,
defaulting to five seconds when no millisecond option is supplied. Thus a
caller's requested duration is not automatically an end-to-end delivery
deadline on this path. Other transports may provide stronger behavior; they
require their own source and runtime review. Do not describe every transport
as asynchronous or assume that a generic timeout parameter bounds all work.
\citep{exec24_openlineage_source}

The Python client adds environment, tags and source-location facets, can
reduce dataset identifiers before filtering, and can return early for absent
transport, disabled emission or a filtered event. Its HTTP transport posts,
checks response status and returns a response, while the public client emit
method returns no response object. The selected retry defaults allow HEAD
and POST and declare total/read/connect counts of five, a backoff factor of
0.3, and status codes 500, 502, 503 and 504. The transport passes this
configuration to urllib3 through Requests. Declared package bounds are not
a solved environment; the effective dependency behavior still needs a locked
profile. A retryable POST and a run UUID do not prove server-side idempotency
or durable exactly-once ingestion. \citep{exec24_openlineage_source}

These distinctions motivate a proposed adapter contract, not an assertion
that it has already been implemented. Preserve a durable local request and
artifact record; emit lineage with a stable logical run identity and explicit
attempt identity; and reconcile the receiver's acknowledgement with the
local record. A response lost after the receiver commits creates ambiguity:
blind repetition can duplicate work unless the receiver's deduplication rule
is known. Cancellation, transport timeout and scientific refusal need separate
fields. Sensitive custom facets also require a logging review because the
inspected clients contain debug event-serialization paths. No actual leakage
or deployment failure is inferred from that observation.
\citep{exec24_openlineage_source}

\subsection{A causal result facet and a request-complete denominator}
Suppose three authored analysis jobs finish their computation, but two
correctly abstain because their identification requirements are unmet. With
a stipulated total cost of USD~45, cost per completed job is USD~15,
whereas cost per accepted causal artifact is USD~45. Reporting only a job
COMPLETE count conceals that distinction. Correct abstention may be valuable;
it should remain explicit rather than disappear from cost or success
denominators. This is a designed accounting example, not a measured service
price or empirical adoption case. \citep{exec24_openlineage_reference}

The proposed facet should therefore state whether the output is an estimate,
identified set, sensitivity analysis, diagnostic, or structured refusal. It
should retain the target and horizon, assumptions and their evidence,
sampling/dependence unit, conditioning population, resolved estimator,
uncertainty object, numerical status, and exact reviewed artifact. An agent
can draft a causal question or graph, but its proposal is not promoted to
verified identification by being placed in an extensible JSON field.
The method owner, artifact reviewer and institutional approver retain their
separate responsibilities.

\subsection{Reference choice, integration alternatives and closure route}
The direct repository license is Apache~2.0. The Java build declares Java~8
source/target with a Java~17 toolchain and a concrete set of library versions;
the Python build declares Python~3.10 or later, bounded build tooling and
minimum runtime dependencies. These declarations support an environment
plan, not transitive rights clearance, a complete lock or a platform support
promise. The preferred conditional route is an original Java or Python client
with a reviewed causal facet and explicit acknowledgement protocol. Direct
schema-compatible emission is an alternative for another language, but it
transfers validation, transport and compatibility maintenance to this project.
A port of every integration is unjustified without a measured requirement.
\citep{exec24_openlineage_source}

For the present product proposal, retain OpenLineage as a versioned protocol
and client reference. Defer native admission until the method owner agrees
the causal vocabulary, the adapter owner freezes clients and receiver,
the rights owner reviews the complete distribution, and the evaluation owner
registers public/generated conformance trials. Trials should cover omitted
states, duplicated terminals, reordered events, suppressed emission, custom
facet rejection, receiver persistence, acknowledgement loss, retry/cancel
behavior and end-to-end costs. A local file lineage backend is an explicit
comparator for small prototypes. Integration with sibling repositories is a
future consumer-specific contract and test, not implied by their presence
on the machine. No sibling repository is modified by this proposal review.

This choice completes a bounded reuse argument while leaving delivery,
native adoption, causal truth and receiver-specific qualification open.
